\subsection{Maximum-Likelihood-Method}
{\scriptsize Funktion $p_\cX(x_1, \ldots; \vartheta) = \P_\vartheta[\cX_1 = x_1, \ldots]$ (gleich wie normale Gew.-F.)}

In Modell $\P_\vartheta$ sind $\overrightarrow{\cX} = (\cX_1, \ldots, \cX_n)$ entweder
\begin{itemize}
    \item \bi{diskret} (gem. Gew. $p_{\overrightarrow{\cX}} (x_1, \ldots, x_n; \vartheta)$)
    \item \bi{stetig} (gem. Dichte $f_{\overrightarrow{\cX}} (x_1, \ldots, x_n; \vartheta)$)
\end{itemize}
Da oft $X_k$ i.i.d mit individueller Gew. $p_\cX(x; \vartheta)$ bzw. Dichte $f_\cX(x; \vartheta)$ sind, also (mit $g$ ersetzt durch $p$ oder $f$)
\[
    g_{\overrightarrow{\cX}}(x_1, \ldots, x_n; \vartheta) = \prod_{k = 1}^n g_\cX(x_k; \vartheta)
\]

\shortdefinition[Likelihood-Funktion]
\[
    L(x_1, \ldots, x_n; \vartheta) = \begin{cases}
        p_{\overrightarrow{\cX}}(x_1, \ldots, x_n; \vartheta) & \text{im diskreten Fall} \\
        f_{\overrightarrow{\cX}}(x_1, \ldots, x_n; \vartheta) & \text{im stetigen Fall}
    \end{cases}
\]
\bi{log-Likelihood-Funktion} $l(\vartheta) = \log(L(x_1, \ldots, x_n; \vartheta))$. Ist im i.i.d-Fall eine Summe.

\inlineintuition Likelihood-F. gibt W., dass in $\P_\vartheta$ die Stichprobe $\cX_1, \ldots, \cX_n$ genau Werte $x_1, \ldots x_n$ liefert.

\shortdefinition[ML-Schätzer] $T_{ML}$ für $\vartheta$ maximiert die Funktion $\vartheta \mapsto L(\cX_1, \ldots, \cX_n; \vartheta)$ für alle $\vartheta$, also
\[
    T_{ML} = t_{TM}(\cX_1, \ldots, \cX_n) \in \underset{\vartheta \in \Theta}{\text{argmax}}\; L(\cX_1, \ldots, \cX_n; \vartheta)
\]
{\scriptsize Wir suchen besten $t_{TM}$, durch geschickte Wahl von $\vartheta$}\\
Meistens sind $\cX_k$ i.i.d. unter $\P_\vartheta$, dann $L$ produkt, also besser $\log(L)$ maximieren (da Summe).

\shortremark
\hl{Einfacher:} Statt max., Nullstellen von Abl. (einfachheitshalber von $l(\vartheta)$) nach $\vartheta$.

\shortremark Suchen oft optimales $L(x_1, \ldots, x_n; \vartheta)$,
da es einfacher ist, ersetzen dann $x_1, \ldots$ durch $\cX_1, \ldots$

\shortexample[Von Serie] Kalibrierung von Gerät bis erfolgreich. $\cX \sim \text{Geo}(\theta)$ Anzahl Versuche, mit $\P_\theta[\cX = k] = \theta (1 - \theta)^{k - 1}$,
gesucht $\theta$, $\cX_i$ unter $\P_\theta$ u.i.v. \bi{Aufgaben}

\underline{LH-Func und $\log$-LH}: Mit Beobachtungen $x_i$ berechnen.\\
$L(x_1, \ldots; \theta) = \ldots = \prod_{k = 1}^n \theta (1 - \theta)^{x_k - 1}$ (s. \ref{sec:series}),\\
$l(\theta) = n \ln(\theta) + (S - n) \ln(1 - \theta)$, mit $S = \sum_{k = 1}^n x_k$.\\
$\log(0)$ möglich (wenn $\theta = 1$ und $S \neq n$, ohne expanding)

\underline{ML-Sch}: $S > n$: $l'(\theta) = \frac{n}{\theta} + \frac{S - n}{1 - \theta}$,
Nullstelle von $l'(\theta)$ ist $\hat{\theta} = \frac{n}{S}$, check mit $l''(\frac{n}{S}) < 0$, ist max. Also ist $T_{ML} = \frac{n}{\sum_{k = 1}^{n} \cX_k}$

\underline{Mit Werten} Einfach einsetzen und auflösen.


% ────────────────────────────────────────────────────────────────────

\shortexample \bi{Verteilungen} (Dimension $\vartheta$ gleich dim. param der V.)\\
\highlight{Bernoulli} $\cX_i \sim \text{Ber}(p)$ i.i.d, hier $\vartheta = p$. Dabei:
$p_\cX(x; \vartheta) = \P_\vartheta[\cX = x] = \vartheta^x (1 - \vartheta)^{1 - x}$ mit $x \in \{0, 1\}$, $S = \sum_{k = 1}^{n}x_k$
\[
    L(x_1, \ldots, x_n; \vartheta) = \vartheta^{S} (1 - \vartheta)^{n - S}
\]
Log LH-Func: $\log(\vartheta) \cdot S + \log(1 - \vartheta) \left( n - S \right)$\\
\bi{ML-Schätzer}: $T = \frac{1}{n} \sum_{k = 1}^{n} \cX_k = \overline{\cX}_n$

\fbox{\bi{Normalverteilung}} $\cX_i \sim \cN(\mu, \sigma^2)$ i.i.d. ($\sigma^2 = v$):
\[
    f_\cX(x; \vartheta) = \frac{1}{\sqrt{2\pi v}} e^{-\frac{(x - \mu)^2}{2v}}
\]
Weil i.i.d (und $\vartheta = (\mu, \sigma^2) = (\mu, v)$):
\[
    L(x_1, \ldots, x_n; \vartheta) = \prod_{k = 1}^n f_\cX(x_k; \vartheta)
\]
und somit ($l(\vartheta) = \log(L(x_1, \ldots, x_n; \vartheta))$):
\[
    l(\vartheta) = -\frac{1}{2} n(\log(2 \pi) + \log(v)) - \sum_{k = 1}^{n} \frac{(x_k - \mu)^2}{2v}
\]
Ableitungen nach $\mu$ und $v$ (beide $0$ für Max, $l$ kurz für $l(\vartheta)$):
\[
    \partial_\mu l = 2 \sum_{k = 1}^{n} \frac{x_k - \mu}{2v} \quad
    \partial_v l = - \frac{n}{2v} + \frac{1}{2v^2} \sum_{k = 1}^n (x_k - \mu)^2
\]

Der \bi{Schätzer} ist $T = (T_1, T_2)$ (\textbf{Momentschätzer}):
\[
    T_1 = \overline{\cX}_n \quad
    T_2 = \frac{1}{n} \sum_{k = 1}^{n} \cX_k^2 - (\overline{\cX}_n)^2
\]
Für $T$ gilt: $(\E_\vartheta[\cX], \V_\vartheta[\cX])$. Ist nicht Erwartungstreu.
Es gilt $\E_\vartheta[\cX_k \cX_l] = \E_\vartheta[\cX_k] \E_\vartheta[\cX_l] = \E_\vartheta[\cX]^2$
und folglich:
\[
    \E_\vartheta[(\overline{\cX}_n)^2] = \frac{1}{n} \E_\vartheta[\cX^2] + \frac{n^2 - n}{n^2} (\E_\vartheta[\cX])^2
\]
\bi{Erwartungstreuer Schätzer} für $(\E_\vartheta[\cX], \V_\vartheta[\cX])$ (s. \ref{sec:estimators}):
\[
    T_1' = T_1 \quad
    T_2' = \frac{n}{n - 1} T_2
\]
{\scriptsize $T_2'$ ist (korrigierte) empirische (Stichproben)varianz \textit{((un)biased sample variance)}.
Bei höherdimensionalen Parametern suchen wir den zugehörigen Nullvektor.}
